%% RSAA THESIS TEMPLATE, BY SIMON MURPHY <simon.murphy.nz@gmail.com>
%% BASED ON THE ORIGINAL TEMPLATE BY LATEX GOD JOSHUA RICH <joshua.rich@gmail.com>
%%
%% README:
% This template is designed to be used with pdflatex rather than plain
% latex.  It was developed under a TeXLive 2010 TeX distribution on
% a Mac, but should also work under Linux and earlier TeX distributions.
%
% For ANU requirements for typesetting and formatting see:
% http://policies.anu.edu.au/guidelines/research_theses_submission_and_examination___information_for_higher_degree_research_students/guideline

%%NOTES ABOUT THE DOCUMENT CLASS
% Yes, I'm using a report style.  Any 'thesis style' you might
% find or someone will give you is based on the standard report
% class, but is probably well outdated compared to whatever TeX
% you have installed.  Besides, why use that style file if you 
% have no idea what it actually does?  You will probably just redefine
% all of its customisations anyway...
\documentclass[a4paper,11pt,openright]{report} 
% ---------------------------------------------------------------------
% command re-definitions and additions
% ---------------------------------------------------------------------

% "Previously published as" chapter starter
\newcommand{\previouslypublished}[1]{\vspace{-0.5cm}\emph{\small #1}}

% i.e. -- e.g. -- etc. -- et. al.
\newcommand{\eg}{{\em e.g.,}}
\newcommand{\ie}{{\em i.e.,}}
\newcommand{\etc}{{\em etc.}}
\newcommand{\etal}{{\em et al.}}
\newcommand{\cf}{{\em cf.}}

% RESOLVE IVOA IDENTIFIERS TO THE US-VO DIRECTORY
% Use this to tag the particular dataset, VO service or tool you used in your research. Users can click on the link and resolve the full description in the US Virtual Observatory directory.
%
% e.g. "I used the NOMAD catalogue, as hosted by the Russian ASTRONET facility (\ivoa{ivo://astronet.ru/cas/nomad})"
%
\newcommand{\ivoa}[1]{\href{http://nvo.stsci.edu/vor10/getRecord.aspx?id=#1}{#1}}

% RESOLVE OBJECTS ON SIMBAD
% Adapted from http://tex.stackexchange.com/questions/4154/how-to-use-newcommand-for-href
% Same functionality as AASTEX macro: http://aastex.aas.org/objects/objectlinking.aas.htm
% #2 is mandatory object name (as written in the text) e.g. \object{IC 2391}
% #1 is optional SIMBAD name (if different from #2) e.g. \object[ETA CHA]{$\eta$~Cha}
\makeatletter
\newcommand{\object}[2][\simbadname]{%
\hypersetup{urlbordercolor=\objectcolor}%
\def\simbadname{#2}%
\StrSubstitute{#1}{ }{+}[\parsedname]%
\href{http://simbad.u-strasbg.fr/simbad/sim-id?Ident=\parsedname}{#2}%
\hypersetup{urlbordercolor=\urlcolor}}%
\makeatother

%% astronomical symbols
\newcommand{\HI}{\hbox{\rmfamily H\,{\textsc i}}}
\newcommand{\HIfat}{\hbox{\rmfamily\bfseries H\,{\textsc i}}}
\newcommand{\HIsub}{\hbox{{\scriptsize H}\,{\tiny I}}}
\newcommand{\HII}{\hbox{\rmfamily H\,{\scshape ii}}\xspace}
\newcommand{\HIIsub}{\hbox{\scriptsize \rmfamily H\,{\scshape ii}}}
\newcommand{\Ha}{\hbox{\rmfamily H\,$\alpha$}}
\newcommand{\msun}{\hbox{$M_{\odot}$}\xspace}
\newcommand{\mhi}{\hbox{$M_{\HIsub}$}}
\newcommand{\lsun}{\hbox{$L_{\odot}$}}
\newcommand{\mlsun}{(M/L)_{\odot}}
\newcommand{\vexp}{\hbox{$V_{exp}$}}
\newcommand{\vhel}{\hbox{$V_{hel}$}}
\newcommand{\vdisp}{\hbox{$\sigma_{disp}$}}
\newcommand{\Nhi}{\hbox{$N_{\HIsub}$}}
\newcommand{\nhi}{\hbox{$n_{\HIsub}$}}
\newcommand{\Mhi}{\hbox{$M_{\HIsub}$}}
\newcommand{\ra}{$\alpha$}
\newcommand{\dec}{$\delta$}
%\newcommand{\degree}{\textdegree}
\newcommand{\arcmin}{\hbox{$^\prime$}}
\newcommand{\arcsec}{\hbox{$^{\prime\prime}$}}
\newcommand{\kms}{\hbox{km s$^{-1}$}}
\newcommand{\mjbeam}{\hbox{mJy beam$^{-1}$}}
\newcommand{\jbeam}{\hbox{Jy beam$^{-1}$}}
\newcommand{\mjbeamkms}{\hbox{mJy/beam km s$^{-1}$}}
\newcommand{\jkms}{\hbox{Jy km s$^{-1}$}}
\newcommand{\coldensity}{\hbox{cm$^{-2}$}}
\newcommand{\voldensity}{\hbox{cm$^{-3}$}}
% add others that you need, or can't be bothered writing repeatedly
\newcommand{\pms}{pre--main sequence}
\newcommand{\Pms}{Pre--main sequence}
\newcommand{\msunyr}{\hbox{$M_{\odot}$~yr$^{-1}$}}
\newcommand{\masyr}{mas~yr$^{-1}$}
\newcommand{\micron}{\hbox{$\mu$m}}
\newcommand{\wifes}{\hbox{\emph{WiFeS}}}
\newcommand{\vsini}{\ensuremath{v\sin i}}

\newcommand{\sprocess}{\textit{s}-process\xspace}
\newcommand{\rprocess}{\textit{r}-process\xspace}

\newcommand{\ud}{\mathrm{d}}
\newcommand{\sun}{\ensuremath{\odot}}
\newcommand{\Mdot}{\ensuremath{\dot{\mathrm{M}}}}
\newcommand{\Msun}{\ensuremath{\mathrm{M}_{\sun}}\xspace}
\newcommand{\Lsub}[1]{\ensuremath{\mathrm{L}_\mathrm{#1}}}
\newcommand{\Lsun}{\ensuremath{{\mathrm L_{\sun}}}}
\newcommand{\Rsun}{\ensuremath{{\mathrm R_{\sun}}}}
\newcommand{\M}[1]{\ensuremath{\mathrm{M}_\mathrm{#1}}}
\newcommand{\Mtwo}[2]{\ensuremath{\mathrm{M}_\mathrm{#1}^\mathrm{#2}}}
\newcommand{\Tsub}[1]{\ensuremath{\mathrm{T}_\mathrm{#1}}}
\newcommand{\Ttwo}[2]{\ensuremath{\mathrm{T}_\mathrm{#1}^\mathrm{#2}}}
\newcommand{\tausub}[1]{\ensuremath{\tau_\mathrm{#1}}}
\newcommand{\Mc}{\M{C}}
\newcommand{\Mcfirst}{\ensuremath{\Mc^\mathrm{TP1}}}
\newcommand{\Mtdu}{\ensuremath{\Mtwo{TDU}{tot}}}
\newcommand{\iso}[2]{\mbox{${}^{#1}{\text{#2}}$}}
%\newcommand{\iso}[2]{\textsuperscript{#1}{\text{#2}}}
\newcommand{\Cratio}{\iso{12}C/\iso{13}C}



% define any objects you will be referring to often
\newcommand{\echa}{\object[ETA CHA CLUSTER]{$\eta$~Cha}} 


% footnote symbols
% use \symbolfootnote[1]{footnote} to get an *
%     * 1 - *
%     * 2 - dagger
%     * 3 - double dagger
%     * 4 - ... 9 (see page 175 of the latex manual) 
\long\def\symbolfootnote[#1]#2{\begingroup%
\def\thefootnote{\fnsymbol{footnote}}\footnote[#1]{#2}\endgroup}
% New definition of square root:
% it renames \sqrt as \oldsqrt
% it defines the new \sqrt in terms of the old one
% See:
%  http://en.wikibooks.org/wiki/LaTeX/Tips_and_Tricks#New_Square_Root
\let\oldsqrt\sqrt
\def\sqrt{\mathpalette\DHLhksqrt}
\def\DHLhksqrt#1#2{%
\setbox0=\hbox{$#1\oldsqrt{#2\,}$}\dimen0=\ht0
\advance\dimen0-0.2\ht0
\setbox2=\hbox{\vrule height\ht0 depth -\dimen0}%
{\box0\lower0.4pt\box2}}

% define a new url command so a nice font can be used
\newcommand{\myurl}[1]{\href{#1}{#1}}

% ---------------------------------------------------------------------
% pdflatex setup
% ---------------------------------------------------------------------
% make pdflatex use the same spacing (paragraph, line and page breaks)
% as standard LaTeX 
\pdfadjustspacing=1

% ---------------------------------------------------------------------
%  misc. options
% ---------------------------------------------------------------------

% ---------------------------------------------------------------------
% more liberal 'float' (tables, figures) placement
% ---------------------------------------------------------------------
% Alter some LaTeX defaults for better treatment of figures:
% See p.105 of "TeX Unbound" for suggested values.
% See pp. 199-200 of Lamport's "LaTeX" book for details.
% General parameters, for ALL pages:
\renewcommand{\topfraction}{0.9}	% max fraction of floats at top
\renewcommand{\bottomfraction}{0.8}	% max fraction of floats at bottom
% Parameters for TEXT pages (not float pages):
\setcounter{topnumber}{2}
\setcounter{bottomnumber}{2}
\setcounter{totalnumber}{4}     % 2 may work better
\setcounter{dbltopnumber}{2}    % for 2-column pages
\renewcommand{\dbltopfraction}{0.9}	% fit big float above 2-col. text
\renewcommand{\textfraction}{0.07}	% allow minimal text w. figs
% Parameters for FLOAT pages (not text pages):
\renewcommand{\floatpagefraction}{0.7}	% require fuller float pages
% N.B.: floatpagefraction MUST be less than topfraction !!
\renewcommand{\dblfloatpagefraction}{0.7}	% require fuller float pages
% remember to use [htp] or [htpb] for placement

% ---------------------------------------------------------------------
% Journal name abbreviations
% ---------------------------------------------------------------------
\newcommand{\jnlref}[1]{\textrm{#1}}
\newcommand{\actaa}{\jnlref{Acta Astronomica}}
\newcommand{\aj}{\jnlref{AJ}}
\newcommand{\araa}{\jnlref{ARA\&A}}
\newcommand{\apj}{\jnlref{ApJ}}
\newcommand{\apjl}{\jnlref{ApJ}}
\newcommand{\apjs}{\jnlref{ApJS}}
\newcommand{\ao}{\jnlref{Appl.~Opt.}}
\newcommand{\apss}{\jnlref{Ap\&SS}}
\newcommand{\aap}{\jnlref{A\&A}}
\newcommand{\aapr}{\jnlref{A\&A~Rev.}}
\newcommand{\aaps}{\jnlref{A\&AS}}
\newcommand{\azh}{\jnlref{AZh}}
\newcommand{\baas}{\jnlref{BAAS}}
\newcommand{\jrasc}{\jnlref{JRASC}}
\newcommand{\memras}{\jnlref{MmRAS}}
\newcommand{\mnras}{\jnlref{MNRAS}}
\newcommand{\pra}{\jnlref{Phys.~Rev.~A}}
\newcommand{\prb}{\jnlref{Phys.~Rev.~B}}
\newcommand{\prc}{\jnlref{Phys.~Rev.~C}}
\newcommand{\prd}{\jnlref{Phys.~Rev.~D}}
\newcommand{\pre}{\jnlref{Phys.~Rev.~E}}
\newcommand{\prl}{\jnlref{Phys.~Rev.~Lett.}}
\newcommand{\pasp}{\jnlref{PASP}}
\newcommand{\pasj}{\jnlref{PASJ}}
\newcommand{\qjras}{\jnlref{QJRAS}}
\newcommand{\skytel}{\jnlref{S\&T}}
\newcommand{\solphys}{\jnlref{Sol.~Phys.}}
\newcommand{\sovast}{\jnlref{Soviet~Ast.}}
\newcommand{\ssr}{\jnlref{Space~Sci.~Rev.}}
\newcommand{\zap}{\jnlref{ZAp}}
\newcommand{\nat}{\jnlref{Nature}}
\newcommand{\iaucirc}{\jnlref{IAU~Circ.}}
\newcommand{\aplett}{\jnlref{Astrophys.~Lett.}}
\newcommand{\apspr}{\jnlref{Astrophys.~Space~Phys.~Res.}}
\newcommand{\bain}{\jnlref{Bull.~Astron.~Inst.~Netherlands}}
\newcommand{\fcp}{\jnlref{Fund.~Cosmic~Phys.}}
\newcommand{\gca}{\jnlref{Geochim.~Cosmochim.~Acta}}
\newcommand{\grl}{\jnlref{Geophys.~Res.~Lett.}}
\newcommand{\jcp}{\jnlref{J.~Chem.~Phys.}}
\newcommand{\jgr}{\jnlref{J.~Geophys.~Res.}}
\newcommand{\jqsrt}{\jnlref{J.~Quant.~Spec.~Radiat.~Transf.}}
\newcommand{\memsai}{\jnlref{Mem.~Soc.~Astron.~Italiana}}
\newcommand{\nphysa}{\jnlref{Nucl.~Phys.~A}}
\newcommand{\physrep}{\jnlref{Phys.~Rep.}}
\newcommand{\physscr}{\jnlref{Phys.~Scr}}
\newcommand{\planss}{\jnlref{Planet.~Space~Sci.}}
\newcommand{\procspie}{\jnlref{Proc.~SPIE}}
\newcommand{\ieeesigprocm}{\jnlref{IEEE~Signal~Processing~Magazine}}
\newcommand{\icarus}{\jnlref{Icarus}}
\let\astap=\aap
\let\apjlett=\apjl
\let\apjsupp=\apjs
\let\applopt=\ao


% define some constants that can be used in various places to easily
% specify title, subjects, author etc.
\newcommand{\thesistitle}{Neutron-Capture Nucleosynthesis and the Chemical Evolution of Globular Clusters}
\newcommand{\fullname}{Luke Jeremy Shingles}
\newcommand{\shortname}{L. J. Shingles}
\newcommand{\thesissubjects}{astronomy,astrophysics}
\newcommand{\thesisdate}{3rd August 2015} %\today
\newcommand{\fullthesisdate}{30th July 2015} %used in disclaimer after signature
\newcommand{\resubmissiondate}{30th September 2015}

\usepackage[british]{babel}
% ---------------------------------------------------------------------
% page size and margins
% ---------------------------------------------------------------------
%%NOTES ABOUT GEOMETRY PACKAGE:
% Below is using the geometry package as it is just cool.  This allows
% you to directly specify the margin requirements as outlined by ANU
% directly.  Here the inner border is 4cm and the outer is 2cm.  Top
% and bottom are also 2cm.  This is much easier than trying to work
% out page and margin dimensions by hand.  This is how you should use
% LaTeX...
\usepackage[pdftex]{geometry}
\geometry{a4paper,twoside,includehead,hmargin={4cm,2.1cm},vmargin={2cm,2cm}}
%\usepackage{showframe} for testing only

% ---------------------------------------------------------------------
% miscellaneous packages used
% ---------------------------------------------------------------------
\usepackage{graphicx}
\usepackage{captcont}
\usepackage{array} % for setting line spacing in tables
\usepackage{epigraph}
\usepackage[nonumberlist,acronym]{glossaries}
\DeclareGraphicsExtensions{.pdf}
%\DeclareGraphicsExtensions{.eps}

\usepackage{import}
\usepackage{verbatim}
\usepackage{listings}
\usepackage{multicol,multirow}
%list of acronyms
\newacronym{AGB}{AGB}{Asymptotic Giant Branch}
\newacronym{BBNS}{BBNS}{Big Bang Nucleosynthesis}
\newacronym{cmd}{CMD}{Colour-Magnitude Diagram}
\newacronym{fdu}{FDU}{First Dredge-Up}
\newacronym{gc}{GC}{Globular Cluster}
\newacronym{hbb}{HBB}{Hot Bottom Burning}
\newacronym{imf}{IMF}{Initial Mass Function}
\newacronym{pdcz}{PDCZ}{Pulse-Driven Convective Zone}
\newacronym{pmz}{PMZ}{Partially Mixed Zone}
\newacronym{pn}{PN(e)}{Planetary Nebula(e)}
\newacronym{RGB}{RGB}{Red Giant Branch}
\newacronym[sort=r-process]{rprocess}{\textit{r}-process}{\textit{rapid} neutron-capture process}
\newacronym{sfr}{SFR}{Star Formation Rate}
\newacronym{sdu}{SDU}{Second Dredge-Up}
\newacronym[sort=s-process]{sprocess}{\textit{s}-process}{\textit{slow} neutron-capture process}
\newacronym{tdu}{TDU}{Third Dredge-Up}
\newacronym{wd}{WD}{White Dwarf}


\usepackage[twoside]{rotating}
\usepackage{ctable} % for good looking tables
\usepackage{longtable}
\usepackage{xstring}
\usepackage{xspace}
\usepackage[bottom]{footmisc}
\usepackage{pdflscape}
% fonts
\usepackage[T1]{fontenc}
\usepackage{textcomp}
\usepackage{ucs}
\usepackage[utf8x]{inputenc}
\usepackage{amsmath,amssymb}
\usepackage{newpxtext}
\usepackage{newpxmath}
%\usepackage{pxfonts}
%\usepackage{tgheros,tgpagella}
\renewcommand*\ttdefault{qcr}
\linespread{1.25} % tgpagella/mathpazo need a bigger line spacing
%\usepackage{times}
%\usepackage{mathptmx} % use times font including math
\usepackage{microtype}

\usepackage{parskip}

%%NOTES ABOUT CAPTION AND SUBFIG PACKAGES
% I'm using the caption package to style my captions.
% You need to read the documentation for this package (on CTAN) as it
% has simple explanations of the options and examples of what it looks
% like.
%
% The subfig package allows the ability to control in detail
% positioning of sub-plots in a complex figure as well as add labels
% to them, which can then be referenced in the text without the need
% for changing the actual number as your sub-figures change.
% The subfig package inherits options from the caption package, so
% declare them together.
%
\usepackage[]{caption,subfig}
\DeclareCaptionLabelSeparator{mysep}{\ \ \ }
\captionsetup{format=plain,%
  indention=0cm,%
  labelsep=mysep,%
  justification=justified,
  labelfont={small,bf},%was footnotesize
  textfont={small}}
% bibliography uses natbib, this will be a requirement
\usepackage[round]{natbib}
% The hyperref package is a must while editing.  It provides urls
% inside your document; you can then click to go to certain pages,
% figures or bibliography entries.
%
% DEFINE MY CUSTOM COLORS
%
% change these to black if you want your citations and links nice and plain
%
\newcommand{\citecolor}{black}
%\newcommand{\citecolor}{violet}
\newcommand{\linkcolor}{black}
\newcommand{\urlcolor}{black}
%\newcommand{\urlcolor}{blue}
\newcommand{\objectcolor}{gray}
%
\usepackage[hyperref,x11names]{xcolor}
\usepackage[unicode,pageanchor,colorlinks=true,plainpages=false,pdfpagelabels,pdftex]{hyperref}
\hypersetup{pdftitle=\thesistitle,%
  pdfauthor=\shortname,%
  pdfsubject={\thesissubjects},%
  citebordercolor=\citecolor,%
  linkbordercolor=\linkcolor,%
  urlbordercolor=\urlcolor,%
  citecolor=\citecolor,%
  linkcolor=\linkcolor,%
  urlcolor=\urlcolor%
}
% Stops figures and tables 'floating' past the section in which they
% are declared.
\usepackage[section]{placeins}
\usepackage[titletoc]{appendix}

\usepackage{ifthen}
\usepackage{xifthen}
% ---------------------------------------------------------------------
% styles for header, footer, pages, bibliography
% ---------------------------------------------------------------------
%

% chapter heading style
% this is the 'Conny' style from the fncychap package, modified to
% remove the two bold lines before the 'Chapter X' title.
% For usage of this package, see:
%  http://tug.ctan.org/cgi-bin/ctanPackageInformation.py?id=fncychap
%
\usepackage{fncychap}
\makeatletter
\ChNameUpperCase
% uncomment line below to convert all your chapter titles to CAPS
%\ChTitleUpperCase  
\ChNameVar{\centering\Huge\usefont{OT1}{qbk}{m}{n}\selectfont}
\ChNumVar{\Huge}
\ChTitleVar{\centering\Huge\usefont{OT1}{qbk}{m}{n}\selectfont}
\ChRuleWidth{2pt}
\renewcommand{\DOCH}{%
  \CNV\FmN{\@chapapp}\space \CNoV\thechapter
  \par\nobreak
  \vskip -0.5\baselineskip
}
\renewcommand{\DOTI}[1]{%
  \mghrulefill{\RW}\par\nobreak
  \CTV\FmTi{#1}\par\nobreak
  \vskip 60\p@
}
\renewcommand{\DOTIS}[1]{%
  \mghrulefill{\RW}\par\nobreak
  \CTV\FmTi{#1}\par\nobreak
  \vskip 60\p@
}

%
%%% page styles:
%
\usepackage{fancyhdr}
\pagestyle{fancy}
\setlength{\headheight}{14pt} 
\renewcommand{\chaptermark}[1]{\markboth{#1}{}} 
\renewcommand{\sectionmark}[1]{\markright{\thesection\hspace{3mm} #1}{}} 
% redefine plain page style
\fancypagestyle{plain}{
\fancyhf{}
\fancyhead[LE]{\usefont{OT1}{qbk}{m}{n}\selectfont \thepage}
\fancyhead[RO]{\usefont{OT1}{qbk}{m}{n}\selectfont \thepage}
\renewcommand{\headrulewidth}{0pt}
\renewcommand{\footrulewidth}{0pt}
}

%
% this next section (till \makeatother) makes sure that blank pages
% are actually completely blank, cause they're not usually
\makeatletter
\def\cleardoublepage{\clearpage\if@twoside \ifodd\c@page\else
	\hbox{}
	\vspace*{\fill}
	\thispagestyle{empty}
	\newpage
	\if@twocolumn\hbox{}\newpage\fi\fi\fi}
\makeatother

%
%%% section header styles:
%
\usepackage[calcwidth]{titlesec}
\titlelabel{\thetitle.\quad}

% uncomment if you would like an index (make sure you actually define
% index terms in your document as well...
%\makeindex
\makeglossaries

% this contains various command definitions and other misc. options
\begin{document}
\pagenumbering{alph} %these are never displayed, but unique page identifiers are needed for hyperref

% epigraph settings, for quotes at the start of introduction and conclusions
\setlength{\epigraphwidth}{10cm}
\setlength{\afterepigraphskip}{1.5cm}
\renewcommand{\epigraphsize}{\small}
\renewcommand{\epigraphrule}{0pt}
\renewcommand{\epigraphflush}{center}
\renewcommand{\textflush}{flushleft}
\citestyle{aa}
\pagestyle{empty}

\newgeometry{twoside,includehead,hmargin={3cm,3cm},vmargin={2cm,2cm}}
%
%%% TITLE PAGE:
%
% This title-page is based on many bits of information found on the
% Internet by searching for 'thesis title page latex'.  Packages are
% also available on CTAN if you don't like this style and don't feel
% confident creating a style yourself.
%
% Remove or change the \usefont..\selectfont if you don't have the
% TeX Gyre Bonum font.
%
\begin{titlepage}\usefont{OT1}{qbk}{m}{n}\selectfont
  \begin{center}
  \phantom{}
  \vspace{0cm}
    \huge{\thesistitle}
    \\[2.5cm]
    \huge{\fullname}
    \\[2.5cm]
    \LARGE{A thesis submitted for the degree of}\\[0.5cm]
      \LARGE{Doctor of Philosophy}\\[0.5cm]
      \LARGE{of the Australian National University}
    \\[3.cm]
     \includegraphics[width=0.5\textwidth]{anu-logo-colour}
     \\[10mm]
        \LARGE{Research School of Astronomy \& Astrophysics}
        \vfill
        \Large{Submitted \thesisdate}\\[1mm]
% Uncomment this line when you print your final library copy        
        \Large{Accepted \resubmissiondate}
  \end{center}
\end{titlepage}
\restoregeometry


%remove if using the parskip package instead
%\setlength{\parindent}{0pt}
\setlength{\parskip}{2.5ex plus 0.5ex minus 0.3ex}

%\cleardoublepage

%\begin{comment}
% PREFACE TO ELECTRONIC COPY
% remember to change bib style file to "astroads" for the PDF copy
%%% COMMENT THIS OUT FOR THE PRINT COPY
%\setcounter{page}{1}
%\phantomsection\addcontentsline{toc}{section}{Notes on the digital copy}
\setlength\fboxsep{3mm}
\setlength\fboxrule{1pt}
\begin{center}
\framebox{\huge{\textsc{Notes on the Digital Copy}}}
\end{center}
\vspace{5mm}
This document makes extensive use of the hyperlinking features of \LaTeX. References to figures, tables, sections, chapters and the literature can be navigated from within the PDF by clicking on the reference. Internet addresses will be displayed in a browser. 

The bibliography at the end of this thesis has been hyperlinked to the NASA Astrophysics Data Service (\myurl{http://www.adsabs.harvard.edu/}). Further information on each of the cited works is available through the link to ADS.

%Several of the large catalogues used in this work have been tagged with their Virtual Observatory unique identifier (e.g. \mbox{ivo://astronet.ru/cas/twomass-psc}). As~various data centres can have different versions and implementations of the same catalogue, this ensures any analysis is repeatable by denoting the exact catalogue \emph{and} service that was used to obtain the data. Clicking on their identifiers will resolve the service descriptions in the  directory hosted by the US Virtual Observatory (\myurl{http://ww.us-vo.org}).

\textbf{This page can be easily removed from the print copy.}

%%% END
%\end{comment}

%
%%% DEDICATION:
%
% dedication page
%\begin{figure}
\centering
{\large\usefont{OT1}{qzc}{m}{n}\selectfont 
To the Australian Taxpayer\dots
\\[4cm]\phantom{1}}
\end{figure}

\cleardoublepage %THIS LINE SWAPS LEFT AND RIGHT PAGES
\normalsize
%
%%% FRONT CONTENT:
%
% page numbering as roman numerals
% don't indent paragraphs 
% increase the amount of space between paragraphs
\pagestyle{plain}
\pagenumbering{roman}
%
% disclaimer
%
% I don't like things showing up in the contents BEFORE the contents
% page itself, but you may wish to uncomment these lines.
%
%\phantomsection\addcontentsline{toc}{section}{Disclaimer}
\section*{Disclaimer}

I hereby declare that, except where acknowledged appropriately either below or in the text, this thesis is my original work undertaken at the Australian National University between February 2012 and August 2015. It has not been submitted for a higher degree at any other university or institution.

This thesis has been submitted as a `Thesis by Compilation' in accordance with the relevant ANU policies. Each of the three main chapters (\ref{chap:shingleskarakas2013}, \ref{chap:shinglesetal2014}, and \ref{chap:shinglesetal2015}) is therefore a completely self-contained article that has been published in a peer-reviewed journal.

Chapter \ref{chap:shingleskarakas2013} has been published in Monthly Notices of the Royal Astronomical Society as \citet*{Shingles:2013kg}. This article includes some results that were previously included in the candidate's Honours thesis in 2011. In 2012, the candidate calculated new nucleosynthesis models to study the impact of the \iso{13}{C} and \iso{22}{Ne} neutron sources and developed the text and figures for publication. Amanda Karakas guided the scientific content of the paper and contributed some sections of text in the Discussion and Conclusions section.

Chapter \ref{chap:shinglesetal2014} has been published in Astrophysical Journal as \citet*{Shingles:2014ja}. The main scientific insights, analysis, figures, and the manuscript were the produced by the candidate with minor revisions in response to feedback from coauthors and the referee. The chemical evolution code used in this project was developed solely by the candidate. The chemical evolution model incorporated tables of yields for massive stars provided by Raphael Hirschi and yields for AGB stars provided by Cherie Fishlock. Machine-readable abundance measurements for M4 and M5 were provided by David Yong.

Chapter \ref{chap:shinglesetal2015} has been published in Monthly Notices of the Royal Astronomical Society as \citet*{Shingles:2015cd}. All evolution and nucleosynthesis models were calculated by the candidate except for the 6 \Msun, $Y=0.40$ nucleosynthesis model, which was calculated by Carolyn Doherty. The figures and all sections of the text were drafted by the candidate and revised with suggestions from coauthors and the referee.

The remainder of this thesis was drafted solely by the candidate and revised in response to comments by Amanda Karakas.

%\vspace{1cm}
\begin{flushright}
  \includegraphics[width=0.3\linewidth]{signature}\\[2mm]
  \fullname\\
  \fullthesisdate
\end{flushright}


\section*{Supplementary Research}
In addition to the work presented in this thesis, the candidate has contributed to several other research projects during the PhD program. A complete list of peer-reviewed journal articles authored or co-authored by the candidate during the 3.5 year program are listed here in reverse chronological order:

\newcommand{\highlightauthor}[1]{\underline{#1}}
\newcommand{\articletitle}[1]{{\flushleft\textbf{#1}}}

\articletitle{\href{http://adsabs.harvard.edu/abs/2015MNRAS.452.2804S}{``Evolution and nucleosynthesis of helium-rich asymptotic giant branch models''}}\\
\highlightauthor{Luke J. Shingles}, Carolyn L. Doherty, Amanda I. Karakas, Richard J. Stancliffe, John C. Lattanzio, Maria Lugaro, 2015, MNRAS, 452, 2804

\articletitle{\href{http://adsabs.harvard.edu/abs/2015MNRAS.450..815M}{``Iron and s-element abundance variations in NGC 5286: comparison with `anomalous' globular clusters and Milky Way satellites''}}\\
A. F. Marino, A. P. Milone, A. I. Karakas, L. Casagrande, D. Yong, \highlightauthor{L. Shingles}, G. Da Costa, J. Norris, P. B. Stetson,  K. Lind, M. Asplund, R. Collet, H. Jerjen, L. Sbordone, A. Aparicio, \& S. Cassisi, 2015, MNRAS, 450, 815

\articletitle{\href{http://adsabs.harvard.edu/abs/2014ApJ...795...34S}{``The s-process enrichment of the globular clusters M4 and M22''}}\\
\highlightauthor{Luke J. Shingles}, Amanda I. Karakas, Raphael Hirschi, Cherie K. Fishlock, David Yong, Gary S. Da Costa, \& Anna F. Marino, 2014, ApJ, 795, 34

\articletitle{\href{http://adsabs.harvard.edu/abs/2014MNRAS.441.3396Y}{``Iron and neutron-capture element abundance variations in the globular cluster M2 (NGC 7089)''}}\\
David Yong, Ian U. Roederer, Frank Grundahl, Gary S. Da Costa, Amanda I. Karakas, John E. Norris, Wako Aoki, Cherie K. Fishlock, A. F. Marino, A. P. Milone, \& \highlightauthor{Luke J. Shingles}, 2014, MNRAS, 441, 3396

\articletitle{\href{http://adsabs.harvard.edu/abs/2013Ap\%26SS.347...47V}{``Augmented reality in astrophysics''}}\\
Fr\'{e}d\'{e}ric Vogt \& \highlightauthor{Luke J. Shingles}, 2013, Ap\&SS, 347, 47

\articletitle{\href{http://adsabs.harvard.edu/abs/2013MNRAS.431.2861S}{``Is the sulphur anomaly in planetary nebulae caused by the s-process?''}}\\
\highlightauthor{Luke J. Shingles} \& Amanda I. Karakas, 2013, MNRAS, 431, 2861

%
% acknowledgements
% \phantomsection\addcontentsline{toc}{section}{Acknowledgements}
\cleardoublepage
\section*{Acknowledgements}
I am thankful for the knowledge, techniques, and culture of science that were developed before my time. Growing up, I had learnt a fair number of scientific facts, but I did so without really believing that these ideas could describe the workings of daily life outside of specialised laboratory experiments. It was far too long before I understood that the known laws of physics are either consistent with literally every aspect of human experience, or we would have evidence that they are wrong and they need to be revised. For me, science then became more important as a way to understand the world while avoiding flawed or irrational thinking.

I am grateful to my primary supervisor, Amanda Karakas for her expert guidance of my research programme. Having observed how heavily other students' progress has been influenced by their supervisors and chosen research topics, I feel very fortunate to have been supervised by Amanda on the topic of stellar evolution and nucleosynthesis. Although she encouraged me to work independently where possible, Amanda has also been very generous with her time and knowledge throughout my PhD. Her excellent teaching skills clearly have applications beyond stellar astrophysics, as she has given me a solid foundation in workplace soft skills such as cocktail making. I'm sure this knowledge will continue to serve me well as I make the transition to postdoctoral life.

Thank you to the rest of my advisory panel, Gary Da Costa, David Yong, Richard Stancliffe, and John Lattanzio. Each of them has given their time to provide me with useful feedback and advice. Thank you also to my collaborators (especially Carolyn Doherty) who have helped me to improve the quality and scientific rigour of my papers.

Thank you to everyone at Stromlo for contributing to a friendly work environment. I would list people by name, but I would inevitably forget someone important, so I will just single out my office mates: Christopher Nolan and Chris Owen. Thanks for waiving the entrance requirements and letting me join the `Chris' office. Thanks to my climbing buddies for helping me to get more exercise and spend less time sitting down. Thanks to everyone at cosmology lunch and supernova tea for welcoming me into the field of supernovae and helping me to prepare for my job interview (and now postdoc) at Queen's University Belfast.

Thank you to my immediate family, Mum, Dad, my sister Tegan, and my brother, Josh. Dr. Joshua Shingles gave me no excuses not to succeed in my doctorate after he completed his while simultaneously working full-time and raising four children with his wife Natalie.

Thank you to Louise for being with me during the best and worst times of our PhD years. I'll remember the great times we spent having tea after work (at exactly 5pm!), spending nights eating in and watching TV at our apartment, having chicken schnitzels and microwaved vegetables ridiculously often, going out for lunch or coffee on the weekends, and our holiday trips together (especially Christmases!). Thank you for introducing me to much of the culture, accents, phrases, and retail stores of the UK ahead of my move to Belfast. I'm reminded of you when I go to a Caff\'e Nero, Boots, Tesco, or Pizza Express, and whenever I have no idea whether a British person is joking or not.
%
% abstract
%\phantomsection\addcontentsline{toc}{section}{Abstract}
\cleardoublepage
\section*{Abstract}

Elements heavier than iron are almost entirely produced in stars through neutron captures and radioactive decays. Of these heavy elements, roughly half are produced by the \gls{sprocess}, which takes place under extended exposure to low neutron densities. Most of the \gls{sprocess} production occurs in stars with initial masses between roughly 0.8 and 8 \Msun, which evolve through the \gls{AGB} phase.

This thesis explores several topics related to \gls{AGB} stars and the \gls{sprocess}, with a focus on comparing theoretical models to observations in the literature on planetary nebulae, post-AGB stars, and globular cluster stars. A recurring theme is the uncertainty of \iso{13}{C}-pocket formation, which is crucial for building accurate models of \sprocess nucleosynthesis.

We first investigated whether neutron-capture reactions in \gls{AGB} stars are the cause of the low sulphur abundances in planetary nebulae and post-AGB stars relative to the interstellar medium. Accounting for uncertainties in the size of the partial mixing zone that forms \iso{13}{C} pockets and the rates of neutron-capture and neutron-producing reactions, our models failed to reproduce the observed levels of sulphur destruction. From this, we concluded that \gls{AGB} nucleosynthesis is not the cause of the sulphur anomaly. We also discovered a new method to constrain the extent of the partial mixing zone using neon abundances in planetary nebulae.

We next aimed to discover the stellar sites of the \gls{sprocess} enrichment in globular clusters that have inter- and intra-cluster variation, with the examples of M4 (relative to M5) and M22, respectively. Using a new chemical evolution code developed by the candidate, we tested models with stellar yields from rotating massive stars and \gls{AGB} stars. We compared our model predictions for the production of \sprocess elements with abundances from s-poor and s-rich populations. We found that rotating massive stars alone do not explain the pattern of abundance variations in either cluster, and that a contribution from \gls{AGB} stars with \iso{13}{C} pockets is required. We derived a minimum enrichment timescale from our best-fitting chemical evolution models and, although the value depends on the assumptions made about the formation of \iso{13}{C} pockets, our estimate of 240--360 Myr for M22 is consistent with the upper limit of 300 Myr inferred by isochrone fitting.

Lastly, there is accumulating evidence that some stars (e.g., in $\omega$ Centauri) have been born with helium mass fractions as high as 40\%. This motivated us to explore the impact of helium-rich abundances on the evolution and nucleosynthesis of intermediate-mass (3--6 \Msun) \gls{AGB} models. We found that the stellar yields of \gls{sprocess} elements are substantially lower in He-rich models, largely as a result of less intershell material being mixed into the envelope. We also found evidence that high He abundances could restrict the \sprocess production by \iso{13}{C} pockets to stars with lower initial masses.

\glsresetall
%
% table of contents, figures and tables
\setcounter{tocdepth}{1} % limit TOC to sections, i.e. 1.2, not the default 1.2.1
%\phantomsection\addcontentsline{toc}{section}{Contents}
\tableofcontents\cleardoublepage
\phantomsection\addcontentsline{toc}{section}{List of Figures}
\listoffigures\cleardoublepage
\phantomsection\addcontentsline{toc}{section}{List of Tables}
\listoftables\cleardoublepage

%
%%% MAIN CONTENT:
%
% page numbering as numbers for content of thesis
\pagenumbering{arabic}
\pagestyle{fancy}
\fancyhead[LO]{\usefont{OT1}{qbk}{m}{n}\selectfont \rightmark}
\fancyhead[RE]{\usefont{OT1}{qbk}{m}{n}\selectfont \leftmark}
\fancyhead[LE]{\usefont{OT1}{qbk}{m}{n}\selectfont \thepage}
\fancyhead[RO]{\usefont{OT1}{qbk}{m}{n}\selectfont \thepage}
\fancyfoot{}
\renewcommand{\headrulewidth}{0.5pt}
\renewcommand{\footrulewidth}{0pt}

\subimport{chapter1/}{chapter1.tex}
\subimport{chapter2/}{chapter2.tex}
\subimport{chapter3/}{chapter3.tex}
\subimport{chapter4/}{chapter4.tex}
\subimport{conclusions/}{conclusions.tex}

%
%%% BACK CONTENT:
%
\cleardoublepage
% for back matter, plain page style with no head/foot rules
%\pagestyle{plain}

% appendices
\cleardoublepage
 \phantomsection
 %
 % I like my appendices to just follow on  from the main text,
 % without a fresh title page
 %
%\appendix
%\appendixpage
%\addappheadtotoc
\begin{appendices}
\subimport{appendices/}{appendices.tex}
\end{appendices}

\phantomsection\addcontentsline{toc}{chapter}{List of acronyms}
\glsaddall
\setglossarystyle{altlist}
\printglossary[type=\acronymtype,title=List of acronyms]

%
% bibliography
\phantomsection\addcontentsline{toc}{chapter}{Bibliography}
% change bib style to 'apj' to remove the ADS links in the printed version
\bibliographystyle{astroads}
%\bibliographystyle{apj}
\bibliography{references}
\end{document}
